\section{Realizing a missing square root}\label{construction}
For the converse to Theorem A we construct an algebra over the given ring from a fraction in its total quotient ring.
\begin{theorem}[Realization over the prescribed base]\label{realize}
Let $(R,\m,k)$ have residue characteristic two. Suppose that either
$2=0$ and $q\in Q(R)\setminus R$ satisfies $q^2\in R$, or
$2$ is a nonzerodivisor and $q\in Q(R)\setminus R$ satisfies $2q,q^2\in R$.
There is a nonsplit module-finite extension $R\subseteq S$ with $S$ local,
$S\otimes_R Q(R)\simeq Q(R)^4$ as a module, and minimal resolution
$0\to R^6\to R^{10}\to S\to0$.
\end{theorem}
\begin{lemma}\label{model}
Take either $P=\mathbf F_2[X,(1+X)^{-1},Y]/(Y^2)$ with $x=X$, or $P=\mathbf Z_{(2)}[Y]/(Y^2)$ with $x=2$. There is a commutative $P$-algebra generated as a module by $1,u,v,w,d$, with relation matrix
\[
 D=\begin{pmatrix}0&x^3Y&0\\0&x^4&0\\Y&0&0\\x^3&Y&0\\0&0&x^2\end{pmatrix}.
\]
It satisfies $x^6(Y+xu)=0$. For any specialization to a local ring $A$ with $x\mapsto a$ regular in its maximal ideal and $Y\mapsto\eta$, the specialized matrix is injective, $A$ embeds in the algebra, and the closed fiber is
\[
 k_A[u,v,w]/(u^3,v^2,w^2,uv,uw,vw).
\]
\end{lemma}
\begin{proof}
Write $M=\operatorname{coker}D$ and $e_0$ for its first generator. The operators $U,V,W$ of Appendix~\ref{cert} induce commuting endomorphisms of $M$. Their values at $e_0$ are $u,v,w$. In the notation of the appendix,
\[
 (U^2-(x^3+h)U-x^2Y\,I)e_0=(1+x)d.
\]
As $1+x$ is invertible in both coefficient rings, $e_0$ is cyclic for the commutative algebra $C=P[U,V,W]$. An operator in $C$ killing $e_0$ kills its entire $C$-orbit, so evaluation at $e_0$ is injective. Hence $C\simeq M$ by evaluation, giving $M$ a unital commutative algebra structure.

The first relation implies $x^3Yw=0$ after multiplication by $Y$, and the second then gives $x^6(Y+xu)=0$. Tensoring with $A$ preserves the presentation. If a linear combination of the three columns has all its entries zero except possibly the first, rows two, four and five force all three coefficients to vanish, in that order, by regularity of $a$. This gives exactness and scalar injectivity. In the closed fiber the matrix vanishes. The operators there give $u^2=d$ and every other product of two generators in $u,v,w,d$ is zero (for products with $d$ use $d=u^2$). This yields the stated local closed fiber.
\end{proof}
\begin{proof}[Proof of Theorem~\ref{realize}]
In equal characteristic write $q=b/a$ with $a$ regular and $b\notin aR$. Set $c=q^2$. Both $a,b$ lie in $\m$, the latter because $b^2=a^2c$. Form
\[
 A=R[t]/(t^2-c),\qquad \eta=b+at .
\]
This is free on $1,t$, and $\eta^2=0$. In mixed characteristic put $a=2$, $b=2q$, $c=q^2$. Now $b^2=4c$, $b\in\m$ and $b\notin 2R$ by regularity of $2$. Put
\[
 A=R[t]/(t^2-bt+c),\qquad \eta=2t-b .
\]
Again $A$ is free on $1,t$ and $\eta^2=b^2-4c=0$. In either case the closed fiber of $A$ is $k[t]/(t^2-\bar c)$, which has one prime. Thus $A$ is local. The element $a$ stays regular and $\eta$ lies in its maximal ideal.

Specialize Lemma~\ref{model} over $A$ and call the resulting algebra $S$. It is finite and local, since its closed fiber over $A$ is local. The maps $R\hookrightarrow A\hookrightarrow S$ are injective. Restriction of the resolution $0\to A^3\to A^5\to S\to0$ gives the indicated ranks over $R$. To see minimality, the only block needing inspection is multiplication by $\eta$ on $A$, whose matrices in the two cases are respectively
\[
 \begin{pmatrix}b&ac\\a&b\end{pmatrix},\qquad
 \begin{pmatrix}-b&-2c\\2&b\end{pmatrix}.
\]
Thus all entries belong to $\m$. Inverting $a$ eliminates $w,u,d$ successively, leaving a free module on $1,v$ over $A[1/a]$. This proves the generic rank assertion.

Finally $a^6(\eta+au)=0$. Every $R$-linear map to $R$ kills $\eta+au$. For a retraction this would give $b+a\rho(t+u)=0$ in equal characteristic, and $2\rho(t+u)-b=0$ in mixed characteristic. Both imply $b\in aR$, a contradiction. Notice that the obstruction is tested over $R$ even though the presentation was constructed over $A$.
\end{proof}
\begin{proof}[Necessity in Theorem A]
If weak normality fails, Lemma~\ref{wn}(3) fails for $p=2$. The fraction witnessing failure satisfies the hypotheses of Theorem~\ref{realize} in each of the two cases of the converse.
\end{proof}
\begin{example}[A cusp and a mixed-characteristic domain]\label{root-examples}
Let $k$ have characteristic two and let
$R=k[[t^2,t^3]]$.  This is the one-dimensional Gorenstein domain
$k[[x,y]]/(y^2-x^3)$.  The fraction $q=t=y/x$ does not belong to
$R$, while $q^2=t^2\in R$.  Theorem~\ref{realize} therefore constructs a
nonsplit rank-four algebra of projective dimension one over this
specific cusp.

For the mixed-characteristic example put
$R=\mathbf Z_2[z]/(z^2-8)$.  It is finite free on $1,z$ over
$\mathbf Z_2$, and its reduction modulo two is
$\mathbf F_2[z]/(z^2)$, so it is complete and local.  The
polynomial is irreducible over $\mathbf Q_2$, since a square in
$\mathbf Q_2$ has even valuation whereas $v_2(8)=3$.  Thus $R$
is a domain.  The element $q=z/2$ lies in its fraction field,
has $q^2=2$ and $2q=z$, and is not in $R$ by the free basis
$1,z$.  The mixed-characteristic part of Theorem~\ref{realize} applies.
\end{example}
\begin{corollary}
A nonreduced Noetherian local ring of characteristic two satisfies $\K$ exactly when its depth is zero. If two is regular and the residue characteristic is two, $\K$ likewise forces reducedness, in every dimension. In each positive-depth case a nonreduced base admits a nonsplit module-finite local extension of generic rank two and Betti numbers $(5,3)$.
\end{corollary}
\begin{proof}
At depth zero, finite projective dimension forces freeness by Auslander--Buchsbaum. In a nonzero finite free algebra the unit is outside $\m S$ by Nakayama, and therefore extends to a basis. At positive depth choose $0\ne z\in R$ with $z^2=0$. In characteristic two take $a$ to be a sufficiently large power of a regular element of $\m$ so that $z\notin aR$, using Krull intersection, and put $\eta=z$. When two is regular take $a=2$ and write $z=2^n\eta$ for $n$ maximal. Then $\eta^2=0$ and $\eta\notin 2R$. Apply Lemma~\ref{model} directly over $R$. Its relation $a^6(\eta+au)=0$ precludes a retraction, and the same elimination and minimality arguments give the assertions.
\end{proof}
